% Generated draft: requires mathematical and source review.
\documentclass[11pt]{article}
\usepackage[margin=25mm]{geometry}
\usepackage{fontspec}
\setmainfont{Latin Modern Roman}
\usepackage{amsmath,amssymb,mathtools,unicode-math}
\setmathfont{Latin Modern Math}
\usepackage{enumitem,needspace,xurl,hyperref}
\setlength{\parindent}{0pt}
\setlength{\parskip}{5pt}
\setlength{\emergencystretch}{3em}
\newcommand{\N}{\mathbb{N}}
\newcommand{\Z}{\mathbb{Z}}
\newcommand{\Q}{\mathbb{Q}}
\newcommand{\R}{\mathbb{R}}
\newcommand{\C}{\mathbb{C}}
\newcommand{\F}{\mathbb{F}}
\newcommand{\A}{\mathbb{A}}
\newcommand{\PP}{\mathbb{P}}
\newcommand{\E}{\mathbb{E}}
\newcommand{\eps}{\varepsilon}
\newcommand{\dd}{\,\mathrm{d}}
\begin{document}
\section*{UnsolvedMath ID: 20003336}
\textbf{Source problem code:} catalogue entry
\subsection{Definitions and mathematical statement}
Question 43 is attributed to A. Katok: do periodic orbits in a triangular billiard correspond to scars, more precisely to quantum limits supported on periodic orbits? The source says that [21] suggested a positive answer for rational billiards, meaning triangles whose angles are rational multiples of $\pi$.

To distinguish the meanings of ``scar'', let $\Omega\subset\R^2$ be a triangle and choose a fixed self-adjoint billiard Laplacian $P$ with specified boundary conditions. For normalized \emph{exact} eigenfunctions
\[
 P u_j=\lambda_j u_j,\qquad \|u_j\|_{L^2(\Omega)}=1,
 \qquad \lambda_j\to\infty,
\]
a quantum limit is a subsequential semiclassical measure on the unit-energy billiard phase space; its position projection is a limit of $|u_j|^2\,dx$. A nonsingular periodic orbit $\gamma$ is a closed billiard trajectory avoiding the vertices. Full support on the orbit means $\mu(\gamma)=1$, which is stronger than a positive periodic component $\mu(\gamma)>0$. The original fragment does not specify Dirichlet versus Neumann boundary conditions, nor its quantifier over triangles and orbits; these must be fixed when applying a precise theorem.

The supplied aid reports that the Hassell--Hillairet--Marzuola vertex-control theorem excludes full support on a nonsingular periodic orbit for true eigenfunctions of polygonal billiards. Its claimed mechanism is a uniform bound $\int_U|u_j|^2\,dx\geq c(U)>0$ in suitably chosen vertex neighborhoods. If the closure $V=\overline U$ avoids the projected orbit, the closed-set limit argument gives $(\pi_*\mu)(V)\geq c(U)$ and hence
\[
 \mu(\gamma)\leq1-c(U),\qquad
 (\pi_*\mu)(O)\leq1-c(U)
\]
for any orbit tube $O$ disjoint from $V$. The constant and admissible boundary-condition hypotheses of the cited theorem are not supplied here. The conclusion excludes a fully orbit-supported quantum limit, not every possible partial scar.

The aid's abstract spectral-transfer obstruction is conditional on the stated vertex bound. Let $E$ be real, $v\in\operatorname{Dom}(P)$ satisfy $\|v\|=1$, and put
\[
 r=\|(P-E)v\|,\qquad \eta=\|1_Uv\|.
\]
Let $\mathcal E$ be a target eigenspace and assume $g=\operatorname{dist}(E,\sigma(P|_{\mathcal E^\perp}))>0$. If every normalized $u\in\mathcal E$ has $\|1_Uu\|^2\geq c>0$, then
\[
 \frac r g\geq\frac{(\sqrt c-\eta)_+}{\sqrt2},
 \qquad t_+=\max(t,0).
\]
Consequently a sequence with $\eta\to0$ and a uniform $c>0$ cannot have $r/g\to0$. This distinguishes localized quasimodes from exact eigenfunctions or a single isolated exact eigenspace.

For the equilateral, right-isosceles, and $30$--$60$--$90$ reflection triangles, the aid reports a stronger torus-unfolding/$L^4$ argument: every quantum limit assigns zero mass to any finite union of individual periodic orbits. That special-family statement must not be generalized to all polygons. The registry says solved and the aid says counterexample/refuted; here these mean the supplied negative full-support conclusion, not an independent certification of all variants of scarring.
\subsection{Short English statement}
Katok asked whether triangular billiards admit exact-eigenfunction quantum limits entirely supported on periodic trajectories. The supplied aid reports a negative answer for nonsingular periodic orbits through vertex control, while distinguishing partial scars and localized quasimodes. It also reports zero periodic-orbit mass for three reflection-triangle families.
\subsection{Sources}
\begin{itemize}
\item Local UnsolvedMath registry, permanent ID 20003336, exact source code catalogue entry; \texttt{aim-other-notes.json}, notes[142], Question 43, A. Katok.
\item American Institute of Mathematics, \emph{Emerging applications of measure rigidity}, \url{https://aimath.org/WWN/measrigid/measrigid.pdf}. The source's 2004-era positive suggestion for rational billiards cites [21]; that bibliography entry is not supplied. The supplied verification says only line-break hyphenation was repaired, but no source visit is claimed here.
\item Supplied research attempt 1: counterexample, problem status refuted, rigor label proved, novelty confidence low; registry status solved and classification SOLVED-IN-LITERATURE. It attributes polygonal vertex control to Hassell--Hillairet--Marzuola and supplies the transfer inequality and special reflection-triangle claim. Machine-generated research aid, not peer reviewed, requiring independent verification before citation. Full publication references and theorem hypotheses are absent.
\end{itemize}
\end{document}
